\chapter{Marco Teórico}

Traducir entre señales cerebrales y espacios latentes visuales exige un piso teórico, y este 
capítulo lo arma. Lo he dividido en cuatro bloques que se apoyan entre sí: la biología de la 
señal neuronal (Neurociencia Computacional), la formalización del mapeo en alta dimensión 
(Matemáticas y Computación), las representaciones semánticas 
(Inteligencia Artificial Multimodal) y los motores que fabrican las imágenes 
(Modelos Generativos Profundos). En cada uno incluyo solo lo justo para sostener las decisiones 
de diseño que aparecen en los capítulos~5 y~6.

\section{Fundamentos Neurocientíficos}

Antes de entender cómo el cerebro interpreta lo que vemos, es importante saber cómo procesa la información visual y cómo se mide esta actividad.

La secuencia que sigo es la siguiente: un estímulo visual llega al cerebro, lo que provoca actividad neuronal; esta actividad neuronal se convierte en una señal BOLD; finalmente, esta señal se traduce en un coeficiente $\beta$ para cada ensayo.

Esta cadena de eventos me permite entender la organización jerárquica del córtex visual. Esta organización es crucial porque justifica por qué divido el cerebro en áreas específicas, llamadas vóxeles, de la manera que lo hago en mi tesis.


\subsection{La Señal BOLD: Acoplamiento Neurovascular y Física de la Medición}

La Resonancia Magnética funcional, o fMRI, no mide directamente la actividad de las neuronas. 
Lo que hace es medir una correlación indirecta con la actividad cerebral. Esto se debe a la conexión 
entre la actividad neuronal y el flujo sanguíneo en el cerebro~\cite{ogawa1990brain,logothetis2001neurophysiological,logothetis2008what}.

Cuando un grupo de neuronas se activa, necesita más oxígeno. El cuerpo responde enviando más sangre 
a esa área. Esto cambia la proporción de oxígeno en la sangre y, como resultado, cambia la señal de 
resonancia magnética. A esta señal se le llama Blood-Oxygen-Level-Dependent, o BOLD.

Sin embargo, la actividad neuronal y la respuesta BOLD no ocurren al mismo tiempo. La respuesta 
BOLD tarda unos 5 o 6 segundos en alcanzar su punto máximo después del estímulo y no vuelve a la 
normalidad hasta después de unos 16 segundos. Esto tiene varias consecuencias importantes para los 
experimentos de fMRI.

En primer lugar, la resolución temporal de la fMRI es de segundos, no de milisegundos como en la 
electrofisiología. En segundo lugar, para analizar los datos de fMRI, se necesitan modelos 
estadísticos bien diseñados, como el esquema Least-Squares-Separate~\cite{mumford2012deconvolving}. Y en tercer lugar, la relación 
señal-ruido en la fMRI es relativamente baja, lo que puede hacer que sea más difícil obtener resultados 
precisos.

La Función de Respuesta Hemodinámica, o HRF, es una herramienta que se utiliza para entender cómo 
la actividad neuronal se relaciona con la respuesta BOLD. El estudio BOLD5000~\cite{chang2019bold5000} utiliza este enfoque 
para analizar los datos de fMRI y obtener coeficientes independientes para cada presentación. Esto 
puede ayudar a mejorar la precisión de los resultados en los experimentos de fMRI.

En BOLD5000 cada vóxel mide unos 2 milímetros en cada dimensión y representa la actividad de cientos 
de miles de neuronas. Entonces, el vector de actividad por estímulo que manejo es el vector 
$\boldsymbol{\beta}_i$ en los números reales con dimensiones $V$, donde $V$ es el número de vóxeles 
dentro de la máscara cortical visual que se ha elegido. 
Dos detalles más dan idea de lo difícil del problema. Uno: la varianza 
entre sujetos ante los mismos estímulos es alta, y de ahí que los modelos por-sujeto sigan 
siendo el estándar en la literatura reciente~\cite{scotti2023mindeye,scotti2024mindeye2,takagi2023,ozcelik2023}. 
Otro: la respuesta BOLD agregada refleja sobre todo potenciales de campo locales de origen 
sináptico, no tanto la salida axonal de las neuronas piramidales~\cite{logothetis2008what}.

\subsection{Jerarquía Cortical Visual y Regiones de Interés (ROIs)}

El córtex visual humano trabaja de forma jerárquica y repartida~\cite{felleman1991distributed,
dicarlo2012does}. La luz que recogen los fotorreceptores de la retina viaja por el nervio óptico 
hasta el núcleo geniculado lateral del tálamo y, de ahí, salta al córtex estriado (V1). Allí 
arranca un procesamiento que pronto se bifurca en dos vías paralelas: la dorsal (el ``dónde'', 
parieto-occipital) y la ventral (el ``qué'', occipito-temporal).Me quedo con la ventral, porque es la que sostiene la representación categórica del contenido visual, que al final es lo que quiero decodificar.
En otras palabras, la ventral es la parte que ayuda a entender las categorías de las imágenes, y eso es lo que me interesa descifrar.

Las áreas tempranas, como V1, V2 y V3, se encargan de codificar características básicas,
como contrastes, bordes orientados, frecuencias espaciales locales y disparidad binocular.
En estas áreas, se pueden ver los principios de selectividad descritos por Hubel y Wiesel.
Cuando la señal llega a V4 y al Complejo Occipital Lateral (LOC), la selectividad se
centra en formas más complejas, partes de objetos y agrupaciones perceptuales.
Más allá de LOC, hay regiones que se especializan en categorías específicas, como
el Área Facial Fusiforme (FFA) para rostros, el Área Parahipocampal de Lugares (PPA)
para escenas y entornos, y el Área Extracorporal (EBA) para partes del cuerpo. Toda esta organización modular está bien cartografiada gracias a estudios de 
codificación retinotópica y categórica~\cite{kay2008identifying,naselaris2011encoding,huth2012continuous,kriegeskorte2008representational}.

De aquí salen tres consecuencias que condicionan el diseño. La primera: la información de una 
imagen no está en un solo vóxel, sino repartida en miles a lo largo de varias ROIs, así que 
hace falta un mapeo $\mathbb{R}^V \rightarrow \mathbb{R}^D$ y no un decodificador vóxel a vóxel. 
La segunda: cada ROI aporta algo distinto ---V1--V3 ponen los rasgos de bajo nivel; V4--LOC, la 
identidad del objeto---, de modo que conviene respetar esa organización al partir el vector de 
vóxeles, cosa que formalizo en el capítulo~5. La tercera: BOLD5000 ya trae una máscara cortical 
visual estándar (V1--V4, LOC, PPA), armada con localizadores retinotópicos y de categoría, que 
baja $V$ de un orden de $\sim 5\!\times\!10^{4}$ a $\sim 6\!\times\!10^{3}$ vóxeles activos 
casi sin perder información discriminativa y, de paso, esquiva el problema combinatorio de la 
sección~\ref{sec:voxel_selection}.

\section{Fundamentos Matemáticos y Computacionales}

Para empezar, ya tenemos la señal BOLD extraída por GLM-LSS. Ahora, el problema es encontrar una forma de decodificarla. Esto se puede ver como un problema de mapeo matemático, que se puede representar como una función $f$ que va de un espacio $\mathbb{R}^V$ a otro espacio $\mathbb{R}^D$. Aquí, $V$ es el número de vóxeles en la máscara cortical y $D$ es la dimensión del espacio latente que queremos alcanzar, que en este caso es $768$, igual que el utilizado en CLIP ViT-L/14.
El tema es que hay muchas más columnas que filas, lo que hace que necesitemos un mapa que sea estable y confiable.
En esta sección se formaliza el régimen $p\!\gg\!n$ propio de la decodificación neural~\cite{hastie2009elements}, 
presenta la regularización de Tikhonov~\cite{tikhonov1977solutions} como solución cerrada 
lineal y se analiza la explosión combinatoria que aparece al querer elegir el mejor subconjunto de 
vóxeles~\cite{garey1979computers}.

\subsection{Alta Dimensionalidad y el Problema $p \gg n$}

Los creadores de los transformadores dicen que estamos ante un caso de lo que ellos llaman la ``maldición de la dimensionalidad''~\cite{hastie2009elements}. Para entender esto mejor, pensemos en una imagen.
Si tenemos que ajustar una recta con mil puntos y solo dos variables, es muy fácil de hacer. Sin embargo, si tenemos que ajustar una recta con cinco mil puntos y cincuenta mil variables, ya no es tan simple. En ese caso, en lugar de ajustar la recta, lo que hacemos es memorizar los datos.

En BOLD5000 los predictores (vóxeles, $V \approx 5\!\times\!10^{3}$--$5\!\times\!10^{4}$ según 
cómo se enmascare) superan por varios órdenes de magnitud a las muestras de entrenamiento 
($N \approx 4\!\times\!10^{3}$ estímulos por sujeto)~\cite{chang2019bold5000}. 
Ese desbalance es el famoso $p \gg n$, y su geometría trae tres problemas. El primero 
es que la matriz de Gram $\mathbf{X}^\top\mathbf{X}\in\mathbb{R}^{V\times V}$ es singular 
(rango $\leq N < V$), con lo cual el sistema normal $\mathbf{X}^\top\mathbf{X}\boldsymbol{\beta} = \mathbf{X}^\top\mathbf{Y}$ 
no tiene solución única: hay infinitas que pasan justo por los datos, y todas fallan fuera de 
la muestra. El segundo es que cualquier estimador por mínimos cuadrados ordinarios se traga el 
ruido térmico y fisiológico del fMRI~\cite{logothetis2008what} y memoriza rarezas en lugar de 
aprender algo que generalice. El tercero es que la varianza del estimador crece con $V/N$ hasta 
dominar al sesgo en el error fuera de muestra el clásico compromiso sesgo-varianza~\cite{hastie2009elements}.

\subsection{Regularización de Tikhonov y Regresión Ridge}

La forma de domar ese problema mal condicionado es la regresión Ridge~\cite{hoerl1970ridge}, 
que en el fondo es lo mismo que la regularización de Tikhonov~\cite{tikhonov1977solutions}. 
Con una matriz de diseño $\mathbf{X}\in\mathbb{R}^{N\times V}$ (la actividad fMRI) y una matriz 
objetivo $\mathbf{Y}\in\mathbb{R}^{N\times D}$ (los embeddings CLIP), basta con sumar una 
penalización $\ell_2$ sobre la magnitud de los coeficientes:
\begin{equation}
\hat{\boldsymbol{\beta}} = \arg\min_{\boldsymbol{\beta}} \|\mathbf{Y} - \mathbf{X}\boldsymbol{\beta}\|_2^2 + \lambda\|\boldsymbol{\beta}\|_2^2
\end{equation}
\myequations{Función de costo de la Regresión Ridge}
Su solución cerrada recupera la invertibilidad con solo sumar una constante $\lambda > 0$ a la diagonal principal:
\begin{equation}
\hat{\boldsymbol{\beta}} = (\mathbf{X}^\top\mathbf{X} + \lambda\mathbf{I})^{-1}\mathbf{X}^\top\mathbf{Y}
\end{equation}
\myequations{Solución cerrada Ridge}
Con una dimensionalidad tan alta, invertir directamente una matriz $V\times V$ sale carísimo. La salida es la formulación dual, apoyada en la Descomposición en Valores Singulares (SVD) $\mathbf{X} = \mathbf{U}\boldsymbol{\Sigma}\mathbf{V}^\top$: baja la complejidad al espacio de las muestras, $\mathcal{O}(N^2 V)$, y deja barrer $\lambda$ sobre una sola descomposición~\cite{hastie2009elements}. Escrita así, la solución dual queda:
\begin{equation}
\hat{\boldsymbol{\beta}}_\lambda = \mathbf{V}\,\text{diag}\!\left(\frac{\sigma_i}{\sigma_i^2 + \lambda}\right)\mathbf{U}^\top\mathbf{Y},
\end{equation}
\myequations{Solución Ridge vía SVD truncado}
y de ahí se lee que cada componente espectral $\sigma_i$ queda atenuada por el factor $\sigma_i^2/(\sigma_i^2+\lambda)$, un filtro pasa-bajos en el espacio singular.

\paragraph{Interpretación Bayesiana.} Hay también una lectura probabilística del estimador 
Ridge. Con verosimilitud Gaussiana 
$\mathbf{Y}\!\mid\!\boldsymbol{\beta} \sim \mathcal{N}(\mathbf{X}\boldsymbol{\beta},\sigma_y^2\mathbf{I})$ 
y prior conjugado $\boldsymbol{\beta}\sim\mathcal{N}(\mathbf{0},\sigma_\beta^2\mathbf{I})$, la 
moda a posteriori (MAP) sale idéntica a la solución cerrada Ridge tomando 
$\lambda = \sigma_y^2/\sigma_\beta^2$. La equivalencia deja claro qué hace el hiperparámetro: 
cuanto más pequeña la varianza esperada de los coeficientes $\sigma_\beta^2$ (prior más fuerte), 
más grande $\lambda$ y más marcada la atenuación espectral.

\paragraph{Condicionamiento numérico y validación cruzada eficiente.} El número de condición 
$\kappa(\mathbf{X}^\top\mathbf{X}+\lambda\mathbf{I}) = (\sigma_1^2+\lambda)/(\sigma_V^2+\lambda)$ 
está acotado por arriba por $\sigma_1^2/\lambda$ cuando $\lambda\!\to\!0^+$, lo que por sí 
solo ya invita a regularizar más allá de lo predictivo. El barrido del hiperparámetro va por 
validación cruzada generalizada (GCV)~\cite{hastie2009elements}: con la SVD hecha una sola vez, 
calcular el error leave-one-out para cualquier $\lambda$ cuesta $\mathcal{O}(N D)$ y no obliga 
a reentrenar.

Todo esto luce impecable sobre el papel. Sin embargo, en la práctica (capítulos~6 y~7) el 
adapter Ridge $\mathbb{R}^V \rightarrow \mathbb{R}^{768}$, ajustado contra los embeddings 
nativos de CLIP ViT-L/14, colapsa en norma, escupe vectores fuera de la distribución de CLIP 
y dispara la ``alucinación tipográfica'' en Stable Diffusion 2.1 unCLIP. El porqué de fondo 
la geometría hiperesférica de CLIP y el Modality Gap lo trato en sección siguiente.

\subsection{Selección de Vóxeles y Complejidad Combinatoria}\label{sec:voxel_selection}

Una idea para bajar la dimensión $p$ sería tirar los vóxeles que aportan poco. El problema es 
que elegir el mejor subconjunto de $k$ vóxeles entre $V$ candidatos hace explotar el conteo 
hasta $\binom{V}{k}$; y si encima se piden restricciones de balanceo entre regiones funcionales 
(V1, V2, V3, V4, LOC, PPA), la tarea se reduce en tiempo polinomial al clásico 3-Partition, 
$NP$-completo en sentido fuerte~\cite{garey1979computers}. Con esa dureza, la búsqueda 
exhaustiva queda descartada.

Quedan al menos tres formas de aproximarse a este problema. Una de ellas es el método greedy, que se aplica a funciones submodulares monótonas. Este método garantiza una aproximación muy buena, de un $(1 - 1/e)$, con respecto a la solución óptima, como se ha demostrado en trabajos anteriores ~\cite{nemhauser1978,krause2014submodular}. Sin embargo, antes de aplicar este método, hay que demostrar que la función que estamos utilizando, ya sea la información mutua o la pérdida del decodificador, es submodular. Esto no es fácil de hacer.

Otra forma de abordar este problema es utilizando la relajación convexa con Lasso o Elastic-Net. Esto convierte el problema de selección discreta en uno continuo. Aunque esto facilita el cálculo, deja algunos coeficientes pequeños que no son nulos y que hay que umbralizar manualmente.

La otra opción que considero es la partición biológica fija. Esta opción sacrifica la optimalidad teórica a cambio de estabilidad y reproducibilidad. Se basa en la idea de que la máscara cortical de BOLD5000, que ha sido validada con localizadores funcionales aparte, contiene casi toda la información discriminativa. Me parece que esta es la opción más eficiente en términos metodológicos. Dejo la extracción de los patrones finos para el adapter no lineal del capítulo 6 y reservo la formulación submodular para la extensión a BCI en tiempo real del capítulo 9.


\section{Fundamentos de Inteligencia Artificial Multimodal}

El “dominio de llegada” de la traducción intermodal es un espacio de representación continuo y amplio llamado espacio latente. En esta parte, voy a revisar las arquitecturas Transformer que lo generan.

Estas arquitecturas utilizan un tipo de aprendizaje llamado aprendizaje contrastivo. Este aprendizaje ayuda a alinear diferentes modalidades dentro del espacio latente.

Una propiedad geométrica importante resulta crucial en este proceso, el Modality Gap, que es lo que empuja a cambiar el adapter 
Ridge por uno contrastivo en el capítulo~6.

\subsection{Mecanismo de Auto-Atención y Vision Transformers}

El Teorema de Aproximación Universal establece que una red neuronal lo suficientemente profunda puede 
aproximar cualquier función continua dentro de un dominio compacto. Los Transformers ~\cite{vaswani2017attention} 
llevan esta idea a la práctica mediante la auto-atención, lo que cambia el enfoque tradicional de los 
sesgos inductivos convolucionales por un cálculo de pesos dinámicos.

Comenzando con una secuencia de tokens, representada como $\mathbf{X}$ en un espacio de 
$\mathbb{R}^{n\times d}$, se generan tres tensores: consultas $\mathbf{Q}$, llaves $\mathbf{K}$ y 
valores $\mathbf{V}$. Estos se calculan mediante la proyección de $\mathbf{X}$ sobre diferentes matrices 
de pesos: $\mathbf{W}_Q$, $\mathbf{W}_K$ y $\mathbf{W}_V$. La atención se calcula utilizando la fórmula:

\begin{equation}
\text{Attn}(\mathbf{Q},\mathbf{K},\mathbf{V}) = \text{softmax}\!\left(\frac{\mathbf{Q}\mathbf{K}^\top}{\sqrt{d_k}}\right)\mathbf{V}.
\end{equation}

Esta fórmula se conoce como atención escalada por producto punto. El factor $1/\sqrt{d_k}$ es importante 
porque evita que el softmax se sature cuando se trabaja con dimensiones altas.

La versión multi-cabeza de los Transformers ejecuta $h$ proyecciones en paralelo. Esto permite capturar 
relaciones en subespacios distintos. Un bloque Transformer completo alterna entre la atención y redes 
eed-forward, conexiones residuales y normalización por capa, según se describe en trabajos como los de 
Ba y Ioffe ~\cite{ba2016layernorm,ioffe2015batchnorm}.


El Vision Transformer (ViT)~\cite{dosovitskiy2021image} lleva esta arquitectura a las imágenes. 
Una imagen $\mathbf{I}\in\mathbb{R}^{H\times W\times 3}$ se trocea en una rejilla de parches 
$14\!\times\!14$ (para ViT-L/14 con entradas $224\!\times\!224$); cada parche se proyecta 
linealmente al espacio de tokens y suma un embedding posicional aprendible. Un token especial 
\texttt{[CLS]} va recogiendo la información global capa a capa y, al final, entrega el embedding 
de la imagen. Como ViT no tiene campos receptivos fijos, arma representaciones contextuales 
densas ya desde la primera capa y se libra de varias limitaciones de localidad propias de las 
convolucionales de siempre.

\subsection{Aprendizaje Contrastivo y CLIP}

CLIP (Contrastive Language--Image Pre-training)~\cite{radford2021learning} entrena a la vez un 
codificador visual ViT y uno textual Transformer sobre 400 millones de pares (imagen, texto) 
sacados de la web. La pérdida es la InfoNCE~\cite{vandenoord2018representation} en su forma 
simétrica bidireccional. Sea un mini-batch de $N$ pares $\{(\mathbf{v}_i,\mathbf{t}_i)\}_{i=1}^{N}$, 
con embeddings visuales $\mathbf{v}_i$ y textuales $\mathbf{t}_i$ llevados al mismo espacio y 
normalizados sobre la hiperesfera unidad; entonces
\begin{equation}
\mathcal{L}_{\text{CLIP}} = -\frac{1}{2N}\sum_{i=1}^{N}\left[\log\frac{e^{\langle\mathbf{v}_i,\mathbf{t}_i\rangle/\tau}}{\sum_{j=1}^{N} e^{\langle\mathbf{v}_i,\mathbf{t}_j\rangle/\tau}} + \log\frac{e^{\langle\mathbf{t}_i,\mathbf{v}_i\rangle/\tau}}{\sum_{j=1}^{N} e^{\langle\mathbf{t}_i,\mathbf{v}_j\rangle/\tau}}\right]
\end{equation}
\myequations{Pérdida InfoNCE simétrica de CLIP}
donde $\tau$ es una temperatura aprendible (por lo común escrita como $\tau = 1/\exp(\texttt{logit\_scale})$, 
con \texttt{logit\_scale} acotada por $\log 100$). Es bidireccional porque alinea a la vez 
imagen$\rightarrow$texto y texto$\rightarrow$imagen; esa misma propiedad la hereda el adapter 
del capítulo~6 para corregir una asimetría empírica entre los gradientes fMRI$\rightarrow$CLIP 
y CLIP$\rightarrow$fMRI.

\subsection{El Fenómeno del Modality Gap}\label{sec:modality_gap}

Lo que más me importa de CLIP no es cómo está construido, sino una peculiaridad geométrica que 
emerge de su entrenamiento: los embeddings visuales y los textuales, aunque se supone que deben 
alinearse, viven en dos conos disjuntos y estrechos de la hiperesfera unidad $\mathbb{S}^{767}\subset\mathbb{R}^{768}$, 
separados por un ángulo casi constante. Liang~et~al.~\cite{liang2022mind} formalizaron el 
fenómeno bautizado "Modality Gap", en línea con observaciones previas de geometría 
representacional~\cite{kriegeskorte2008representational}. Sus tres consecuencias para el 
problema fMRI$\rightarrow$CLIP son las que justifican el giro arquitectónico del capítulo~6.

Va la primera: el espacio CLIP no es uniformemente denso, hasta el punto de que una gaussiana 
isotrópica en $\mathbb{R}^{768}$ le daría medida cero al cono visual real. Un predictor fMRI
$\rightarrow$CLIP entrenado con MSE puro tiende a la media empírica del espacio, cae fuera del 
cono y termina pasándole al decodificador entradas fuera de distribución (OOD). La segunda 
tiene que ver con el propio Ridge: óptimo en $\ell_2$, pero ciego a la dirección hiperesférica, 
encoge la norma de las predicciones hacia valores intermedios y empeora el OOD, porque ahora 
las predicciones se alejan de la distribución real de embeddings tanto en dirección como en 
magnitud. Y la tercera es la que cierra el círculo: Stable Diffusion 2.1 unCLIP solo vio 
embeddings CLIP nativos durante su entrenamiento, así que al recibir un vector OOD lo lee como 
una orden ambigua y se refugia en su prior más fuerte, el de los tokens tipográficos, dibujando 
texto en vez de objetos coherentes. Es la ``alucinación tipográfica'' que documento en los 
capítulos~7 y~8.

De ahí sale una conclusión metodológica directa: para acortar el Modality Gap, aunque sea en 
parte, hace falta un objetivo sensible al ángulo y no a la distancia euclídea. Lo natural sería 
sustituir el Ridge por un adapter no lineal contrastivo (InfoNCE bidireccional, con MSE y 
coseno como anclas auxiliares), que es la receta con la que la familia MindEye~\cite{
scotti2023mindeye,scotti2024mindeye2} marcó el estado del arte y que dejo como continuación en 
el Cap.~9. Pero antes de pagar ese precio quiero saber cuánto del fallo del Ridge se debe solo 
a que el embedding llega sin variabilidad ni norma al decodificador. Para medirlo uso el Ridge 
Estocástico, que formalizo en la sección siguiente.

\subsection{Insuficiencia del MSE y Justificación del Objetivo Contrastivo}

Tomemos un par $(\mathbf{f},\mathbf{c}^{*})$, donde $\mathbf{f}\in\mathbb{R}^{D}$ es lo que 
predice el adapter y $\mathbf{c}^{*}\in\mathbb{R}^{D}$ el embedding CLIP objetivo. La pérdida 
MSE
\begin{equation}
\mathcal{L}_{\text{MSE}}(\mathbf{f},\mathbf{c}^{*}) = \|\mathbf{f}-\mathbf{c}^{*}\|_2^2 = \|\mathbf{f}\|_2^2 - 2\langle\mathbf{f},\mathbf{c}^{*}\rangle + \|\mathbf{c}^{*}\|_2^2
\end{equation}
\myequations{Descomposición del MSE en magnitud y dirección}
La combinación de dos objetivos que no son compatibles es un problema. Por un lado, se intenta reducir la 
diferencia entre la magnitud de $\mathbf{f}$ y la magnitud de $\mathbf{c}^{*}$, es decir, 
$\|\mathbf{f}\|_2^2 - \|\mathbf{c}^{*}\|_2^2$. Por otro lado, se busca aumentar el producto interno 
$\langle\mathbf{f},\mathbf{c}^{*}\rangle$.

En un espacio como el de CLIP, donde el significado se representa mediante ángulos y no mediante normas, 
el primer término introduce un gradiente que no es deseado. Esto hace que todas las predicciones se muevan 
hacia la norma media del conjunto de datos. Este es el colapso que se observa cuando se utiliza Ridge.

Por otro lado, la función de pérdida InfoNCE juega un papel diferente. Trabaja con productos internos 
normalizados, como $\hat{\mathbf{f}}\cdot\hat{\mathbf{c}}^{*} = \cos\theta$, y busca maximizar la 
información mutua entre el par positivo $(\hat{\mathbf{f}}_i,\hat{\mathbf{c}}^{*}_i)$ y los pares 
negativos $\{(\hat{\mathbf{f}}_i,\hat{\mathbf{c}}^{*}_j)\}_{j\neq i}$ dentro del mismo lote de 
entrenamiento~\cite{vandenoord2018representation}. Su objetivo es
\begin{equation}
\mathcal{L}_{\text{NCE}}^{f\rightarrow c} = -\frac{1}{N}\sum_{i=1}^{N}\log\frac{\exp(\hat{\mathbf{f}}_i\cdot\hat{\mathbf{c}}^{*}_i/\tau)}{\sum_{j=1}^{N}\exp(\hat{\mathbf{f}}_i\cdot\hat{\mathbf{c}}^{*}_j/\tau)},
\end{equation}
\myequations{InfoNCE direccional fMRI$\rightarrow$CLIP}
y su versión bidireccional añade el término simétrico 
$\mathcal{L}_{\text{NCE}}^{c\rightarrow f}$. 
El gradiente que sale es puramente angular y no tira de la norma de $\mathbf{f}$ hacia ningún lado, por lo que el colapso de norma del Ridge 
desaparece. La receta híbrida de MindEye~\cite{scotti2023mindeye} combina InfoNCE bidireccional con un 
peso de 1.0, el error cuadrático medio con un peso de 0.3 y el coseno con un peso de 0.3.
En este caso, los dos términos auxiliares actúan como anclas geométricas suaves sin robarle protagonismo 
al gradiente principal. No entreno esa pérdida, en su lugar, aplico la corrección estocástica de norma y 
dirección que describo en la sección ~\ref{sec:ridge_stoch}.

\subsection{Bloques Residuales (referencia para la línea contrastiva)}

Los adapters contrastivos profundos de MindEye se apoyan en dos piezas clásicas del deep 
learning moderno: bloques residuales con conexiones de salto~\cite{he2016deep} y normalización 
por capa en pre-activación~\cite{ba2016layernorm}. Un bloque residual hace
\begin{equation}
\mathbf{h}_{l+1} = \mathbf{h}_l + \text{MLP}\!\left(\text{LayerNorm}(\mathbf{h}_l)\right),
\end{equation}
\myequations{Bloque residual pre-norm}
donde $\text{MLP}(\cdot)$ es una red feed-forward de dos capas con activación GELU. La 
conexión $+\mathbf{h}_l$ deja pasar el gradiente sin que se atenúe, y así resuelve el problema 
el vanishing gradients de las MLPs profundas. Se prefiere la normalización por ~\cite{ioffe2015batchnorm} 
capa en lugar de la normalización por lote porque cuando se trabaja con 
lotes pequeños, como los de BOLD5000, las estadísticas de cada lote dejan de ser fiables.
Por ejemplo, la configuración typical de MindEye utiliza cuatro de estos bloques con una dimensión 
oculta de 4096, lo que resulta en unos 154 millones de parámetros que se pueden entrenar.
No entreno esta arquitectura en la tesis, pero la documento porque es una extensión natural del trabajo 
que se presenta en el Capítulo 9 y porque compararla con el Ridge Estocástico permite medir el costo y el 
beneficio de dar el salto a un enfoque contrastivo completo.

\section{Ridge Estocástico: Formalización Teórica}\label{sec:ridge_stoch}

El Ridge Estocástico es la base de esta tesis. Corrige de manera sencilla el Ridge Lineal.
Devuelve dos características geométricas del manifold CLIP: la norma estándar y la
variabilidad direccional. Esto se logra sin necesidad de entrenar un adaptador no lineal.
La idea es simple: el Ridge siempre proporciona el mismo vector promedio y aplanado.
Solo hay que moverlo un poco y volver a unirlo a la esfera donde se encuentran los
embeddings reales.La idea la tomo de la familia MindEye~\cite{scotti2023mindeye,scotti2024mindeye2}. 
Ellos mostraron que la geometría angular de CLIP~\cite{liang2022mind} es lo que realmente importa al 
decodificar. También me basé en la idea de que la regularización estocástica puede verse como un 
suavizado del estimador~\cite{hastie2009elements}. A continuación, formalizaré la construcción de esta
idea. También explicaré por qué corrige el fallo del método Ridge. Además, mostraré cómo el ruido $\sigma$ 
se relaciona con la métrica de evaluación final.


\subsection{Definición y motivación geométrica}

Llamemos $\hat{\boldsymbol{\beta}}\in\mathbb{R}^{V\times D}$ a la solución cerrada del Ridge 
Lineal~\cite{hoerl1970ridge,tikhonov1977solutions} y $\hat{\mathbf{e}}_{\text{Ridge}} = \mathbf{X}\hat{\boldsymbol{\beta}}$ 
a la predicción de un trial. El Ridge Estocástico la retoca en dos pasos:
\begin{equation}
\hat{\mathbf{e}}_{\text{stoch}} = \hat{\mathbf{e}}_{\text{Ridge}} + \sigma\,\boldsymbol{\xi}, \quad \boldsymbol{\xi}\sim\mathcal{N}(\mathbf{0},\mathbf{I}_D),
\end{equation}
\myequations{Predictor Ridge Estocástico}
\begin{equation}
\hat{\mathbf{u}}_{\text{stoch}} = \frac{\hat{\mathbf{e}}_{\text{stoch}}}{\|\hat{\mathbf{e}}_{\text{stoch}}\|_2} \in \mathbb{S}^{D-1}.
\end{equation}
\myequations{Proyección sobre la hiperesfera unidad}
Traducido: Primero tomo la salida del Ridge y la mezclo con ruido blanco, luego la adjunto a la esfera 
unidad. El ruido Gaussiano se distribuye de manera uniforme en todas direcciones, pero cuando ajusto la 
densidad, se concentra en un cono dentro de la esfera. La dirección del eje de este cono es la misma que 
la de $\hat{\mathbf{e}}_{\text{Ridge}}$. La abertura del cono aumenta a medida que $\sigma$ se vuelve más 
grande en comparación con $\|\hat{\mathbf{e}}_{\text{Ridge}}\|$. Esto no es algo nuevo en el ámbito de 
$\mathbb{S}^{D-1}$. Técnicas como el aprendizaje contrastivo ya utilizan esta geometría para asegurar 
que las representaciones sean alineadas y uniformes~\cite{vandenoord2018representation,liang2022mind}, 
según la cual, los embeddings de clases similares deben tener la misma dirección, pero no necesariamente 
la misma magnitud.
Con $\sigma$ mucho menor que la norma de $\hat{\mathbf{e}}_{\text{Ridge}}$, obtengo una dirección muy 
parecida a la de Ridge, pero renormalizada.
En cambio, cuando $\sigma$ es mucho mayor que la norma de $\hat{\mathbf{e}}_{\text{Ridge}}$, se obtiene 
algo cercano a una distribución uniforme sobre la esfera $\mathbb{S}^{D-1}$.
El objetivo es encontrar un punto intermedio que maximice la separabilidad direccional entre clases 
visuales sin que se pierdan en el ruido.
Y ese punto es precisamente el que se busca en el barrido de $\sigma$ que se presenta en el capítulo~6.

\subsection{Lectura como tempering posterior}
En el modelo bayesiano de Ridge, se utiliza el modelo gaussiano para obtener una estimación de los coeficientes, que se denota por $\hat{\boldsymbol{\beta}}_{\text{MAP}}$.

Este modelo se basa en el trabajo de Hoerl y Hastie, quienes lo describieron en sus artículos \cite{hoerl1970ridge,hastie2009elements}.

Cuando tomamos una muestra de $\boldsymbol{\xi}$, es similar a tomar una muestra de la distribución del modelo.

Esta distribución se puede escribir como:

\begin{equation}
p(\mathbf{e}\!\mid\!\mathbf{x},\mathcal{D}) \approx \mathcal{N}\!\left(\mathbf{X}\hat{\boldsymbol{\beta}}_{\text{MAP}},\,\Sigma_{\text{pred}}\right)
\end{equation}
\myequations{Aproximación isotrópica del posterior predictivo}
que es una distribución normal, con una media de $\mathbf{X}\hat{\boldsymbol{\beta}}_{\text{MAP}}$ y una covarianza de $\Sigma_{\text{pred}}$.

El modelo bayesiano de Ridge se utiliza para realizar predicciones y es un tipo de modelo de regresión que se emplea para estimar los coeficientes, es decir, $\hat{\boldsymbol{\beta}}_{\text{MAP}}$.

con $\Sigma_{\text{pred}}$ aproximada por la covarianza isotrópica $\sigma^2\mathbf{I}_D$. 
Visto así, el modelo deja de devolver siempre la misma media: admite que no está seguro y 
muestrea a su alrededor. El papel de la temperatura $\sigma$ es equilibrar exploración y 
explotación, una idea que viene de los generativos basados en denoising score matching, donde 
la varianza del ruido dosifica la mezcla entre estimador puntual y diversidad
~\cite{sohldickstein2015deep,ho2020ddpm}. Con $\sigma=0$ recupero la inferencia MAP pura, es 
decir el Ridge Lineal; con $\sigma>0$ inyecto incertidumbre epistémica y saco la predicción de 
las zonas OOD donde el decodificador se cae.

\subsection{Por qué la renormalización es necesaria}

Si me saltara la proyección esférica, el ruido Gaussiano cambiaría la norma esperada de la 
predicción de un modo totalmente previsible: $\mathbb{E}\|\hat{\mathbf{e}}_{\text{stoch}}\|_2^2 = \|\hat{\mathbf{e}}_{\text{Ridge}}\|_2^2 + D\sigma^2$, con lo cual solo empeoraría el problema de norma OOD. Renormalizar con $\hat{\mathbf{u}}_{\text{stoch}} = \hat{\mathbf{e}}_{\text{stoch}}/\|\hat{\mathbf{e}}_{\text{stoch}}\|_2$ fija la magnitud en $1$ ---la norma típica de los embeddings CLIP nativos~\cite{liang2022mind}--- y separa lo que el ruido hace en dirección de lo que haría en magnitud, dejando solo su aporte angular.

\subsection{Calibración de $\sigma$ y pairwise accuracy}

Para calibrar $\sigma$ no minimizo el MSE ---que, como vimos en~\ref{sec:modality_gap}, 
es ciego a la dirección---, sino que maximizo de una vez la métrica final: la pairwise accuracy 
en sentido 2-AFC sobre embeddings CLIP de validación, hoy el criterio habitual en la 
decodificación visual~\cite{scotti2023mindeye,scotti2024mindeye2,ozcelik2023,takagi2023}. 
Dicho de otra manera: en vez de pedirle al modelo que clave el vector exacto, le pido algo más 
flexible y a la vez más provechoso ---que el embedding predicho se parezca más al de la imagen 
vista que al de cualquier otra del conjunto---. En concreto, para cada par 
$(\mathbf{e}_i,\mathbf{e}_j)$ con $i\!\neq\!j$ mido
\begin{equation}
\mathbb{P}\!\left[\langle\hat{\mathbf{u}}_i^{\text{stoch}},\mathbf{e}_i\rangle > \langle\hat{\mathbf{u}}_i^{\text{stoch}},\mathbf{e}_j\rangle\right],
\end{equation}
\myequations{Pairwise accuracy en CLIP}
y me quedo con $\sigma^{*} = \arg\max_\sigma$ de esa probabilidad estimada sobre validación. 
Así el objetivo de calibración queda pegado al criterio con el que luego comparo contra el SOTA 
en el capítulo~7~\cite{scotti2023mindeye}, y de paso evito el clásico desajuste entre la 
métrica con la que se entrena y la que se reporta~\cite{hastie2009elements}.

\section{Modelos Generativos}

Falta un último eslabón: la representación latente que sale del adapter tiene que volver a 
píxeles, y de eso se encarga un decodificador probabilístico expresivo. Lo veo como un reparto 
de tareas ---el adapter marca un ``rumbo'' en el espacio semántico y el decodificador dibuja la 
imagen que le corresponde---. Repaso aquí los modelos de difusión latente
~\cite{ho2020ddpm,sohldickstein2015deep,rombach2022high}, el condicionamiento por embedding 
visual de unCLIP~\cite{ramesh2022unclip} y la Classifier-Free Guidance
~\cite{ho2022cfg,ho2022classifier}, con la que refuerzo la fidelidad semántica de las 
reconstrucciones.

\subsection{Modelos de Difusión: Proceso Directo y Reverso}

Los modelos de difusión~\cite{sohldickstein2015deep,ho2020ddpm} son procesos generativos 
estocásticos: aprenden a convertir ruido gaussiano puro en muestras estructuradas siguiendo una 
cadena de Markov que avanza a pasos pequeños. El proceso directo va ensuciando un dato 
$\mathbf{z}_0$ poco a poco, en $T$ pasos discretos, según
\begin{equation}
q(\mathbf{z}_t\mid\mathbf{z}_{t-1}) = \mathcal{N}\!\left(\mathbf{z}_t;\,\sqrt{1-\beta_t}\,\mathbf{z}_{t-1},\,\beta_t\mathbf{I}\right),
\end{equation}
\myequations{Proceso directo de difusión}
donde el cronograma $\{\beta_t\}_{t=1}^{T}$ fija cuánta varianza de ruido entra en cada paso. 
Un detalle analítico clave es que existe la marginal cerrada 
$q(\mathbf{z}_t\mid\mathbf{z}_0) = \mathcal{N}(\sqrt{\bar{\alpha}_t}\,\mathbf{z}_0,\,(1-\bar{\alpha}_t)\mathbf{I})$ 
con $\bar{\alpha}_t = \prod_{s\leq t}(1-\beta_s)$, que deja muestrear $\mathbf{z}_t$ 
directamente desde $\mathbf{z}_0$ sin recorrer toda la cadena.

El proceso reverso pone una red neuronal $\epsilon_\theta(\mathbf{z}_t,t,\mathbf{c})$,
casi siempre una U-Net con bloques de atención cruzada que, condicionada por un contexto 
$\mathbf{c}$, aprende a estimar el ruido que se inyectó en cada paso. Su objetivo de 
entrenamiento es
\begin{equation}
\mathcal{L}_{\text{LDM}} = \mathbb{E}_{\mathbf{z}_0,\boldsymbol{\epsilon},t,\mathbf{c}}\!\left[\|\boldsymbol{\epsilon} - \epsilon_\theta(\mathbf{z}_t,t,\mathbf{c})\|_2^2\right].
\end{equation}
\myequations{Pérdida condicional de los modelos de difusión}
Los Latent Diffusion Models (LDM)~\cite{rombach2022high,rombach2022ldm} corren todo esto no en 
píxeles, sino en el espacio latente comprimido $\mathbf{z}\in\mathbb{R}^{4\times h/8\times w/8}$ 
de un autoencoder VAE preentrenado. El truco recorta el costo en dos órdenes de magnitud sin 
que la calidad perceptual se resienta.

\subsection{unCLIP: Condicionamiento por Embedding Visual}

Stable Diffusion 2.1 unCLIP~\cite{ramesh2022unclip,rombach2022high} es una variante del LDM que 
cambia el contexto textual de siempre (el que sale de un encoder CLIP de texto) por un embedding 
visual CLIP ViT-L/14, inyectado directamente en el parámetro \texttt{image\_embeds} de la pipeline.
A nivel de arquitectura, en lugar del encoder textual congelado hay una proyección lineal que 
lleva el embedding visual al espacio de atención cruzada de la U-Net. Para esta tesis la ventaja 
es doble y evidente: cualquier vector $\mathbf{f}\in\mathbb{R}^{768}$ que produzca el adapter 
fMRI$\rightarrow$CLIP entra como condicionamiento sin pasar por el encoder de texto. Así se 
cortan las fugas semánticas por el prompt positivo y queda garantizado que las alucinaciones 
tipográficas del Ridge vienen del embedding del adapter, no del encoder textual.

Esta decisión también sirve como herramienta de diagnóstico pues al dejar el "prompt" de texto 
en cadena vacía y CFG opera con embedding nulo $\emptyset$, todo sesgo categórico que 
aparezca en las reconstrucciones (texto, rostros, escenas) sale solo del embedding predicho. 
Eso me permite leer el modo de fallo como un termómetro del estado del adapter.

\subsection{Acondicionamiento por Cross-Attention y Classifier-Free Guidance}

El vector $\mathbf{c}$ entra a la U-Net por atención cruzada en varias resoluciones: las 
consultas salen de la representación ruidosa $\mathbf{z}_t$ y las llaves/valores de la 
proyección de $\mathbf{c}$, de modo que en cada bloque se mezclan lo espacial y lo semántico. 
Para pegar más la imagen al vector de condicionamiento uso Classifier-Free Guidance (CFG)
~\cite{ho2022cfg,ho2022classifier}, que extrapola sobre la marcha entre la predicción 
condicionada $\epsilon_\theta(\mathbf{z}_t,\mathbf{c})$ y la incondicionada 
$\epsilon_\theta(\mathbf{z}_t,\emptyset)$:
\begin{equation}
\tilde{\epsilon}_\theta(\mathbf{z}_t,\mathbf{c}) = \epsilon_\theta(\mathbf{z}_t,\emptyset) + w\!\cdot\!\left(\epsilon_\theta(\mathbf{z}_t,\mathbf{c}) - \epsilon_\theta(\mathbf{z}_t,\emptyset)\right),
\end{equation}
\myequations{Classifier-Free Guidance}
donde $w>1$ es la escala de guía. En mi caso los valores útiles caen en $w\in[5,9]$: subirlo 
aprieta la fidelidad semántica al embedding a costa de diversidad y, curiosamente, también 
intensifica las alucinaciones tipográficas cuando el embedding está fuera de distribución. 
Ese detalle cierra el lazo diagnóstico que abrí en la sección~\ref{sec:modality_gap}: el 
problema del Ridge no se arregla subiendo $w$, sino corrigiendo la geometría del embedding que 
se predice. Y esa es, en el fondo, la razón metodológica para pasar al Ridge Estocástico que 
formalizo en la sección~\ref{sec:ridge_stoch} y pongo a prueba en los capítulos~6 y~7.

\afterpage{\blankpage}
